% GENERATED FILE. Edit canonical lessons, then run npm run generate:books.
% canonical-source-hash: fnv1a64:7336733fd281292e
% canonical-lessons: SA-C192-harah, SA-C192-mukutam, SA-C192-adhovastram, SA-C192-paridhanam, SA-C192-vasanam

\chapter{Necklace, Crown, Lower garment, Clothing, Garment}
\label{ch:sa-necklace-crown-lower-garment-clothing-garment}
\hlchaptermodality{192}

\begin{chapteropening}
\textbf{By the end of this chapter:} \emph{I can say a necklace, a crown, a lower garment, clothing and a garment.}

Complete the last lesson of chapter 192: I can say a necklace, a crown, a lower garment, clothing and a garment.
\end{chapteropening}

\section[\texorpdfstring{\sk{हारः} (hāraḥ)}{हारः (hāraḥ)}]{\sk{हारः} (hāraḥ) — a necklace}

\label{lesson:SA-C192-harah}



\begin{quote}
Before the new one: say the Sanskrit for a pestle, then the Sanskrit for a sieve.
\end{quote}

\subsection*{You'll want to know: \sk{हारः}}
\textbf{\sk{हारः}} — \emph{hāraḥ} — \textquotedblleft{}a necklace\textquotedblright{}.

\begin{grammarlens}[title={What you've built}]
One more word for this chapter.
\end{grammarlens}

\subsection*{Your turn}
\begin{itemize}
\raggedright
  \item \emph{Say it:} \emph{hāraḥ}
  \item \emph{Say it:} \emph{hāraḥ}, once more
  \item \emph{Say it:} say \emph{hāraḥ}
  \item \emph{Recall:} say the Sanskrit for a pestle, then the Sanskrit for a sieve, then say \emph{hāraḥ} again
\end{itemize}

\begin{tcolorbox}[breakable,colback=teal!4,colframe=teal!35!black,title={Before you move on}]
What does \sk{हारः} mean? (\textquotedblleft{}A necklace\textquotedblright{}.) Say it once more.
\end{tcolorbox}
\section[\texorpdfstring{\sk{मुकुटम्} (mukuṭam)}{मुकुटम् (mukuṭam)}]{\sk{मुकुटम्} (mukuṭam) — a crown}

\label{lesson:SA-C192-mukutam}



\begin{quote}
Before the new one: say the Sanskrit for a sieve, then the Sanskrit for a necklace.
\end{quote}

\subsection*{You'll want to know: \sk{मुकुटम्}}
\textbf{\sk{मुकुटम्}} — \emph{mukuṭam} — \textquotedblleft{}a crown\textquotedblright{}.

\begin{grammarlens}[title={What you've built}]
One more word for this chapter.
\end{grammarlens}

\subsection*{Your turn}
\begin{itemize}
\raggedright
  \item \emph{Say it:} \emph{mukuṭam}
  \item \emph{Say it:} \emph{mukuṭam}, once more
  \item \emph{Say it:} say \emph{mukuṭam}
  \item \emph{Recall:} say the Sanskrit for a sieve, then the Sanskrit for a necklace, then say \emph{mukuṭam} again
\end{itemize}

\begin{tcolorbox}[breakable,colback=teal!4,colframe=teal!35!black,title={Before you move on}]
What does \sk{मुकुटम्} mean? (\textquotedblleft{}A crown\textquotedblright{}.) Say it once more.
\end{tcolorbox}
\section[\texorpdfstring{\sk{अधोवस्त्रम्} (adhovastram)}{अधोवस्त्रम् (adhovastram)}]{\sk{अधोवस्त्रम्} (adhovastram) — a lower garment}

\label{lesson:SA-C192-adhovastram}



\begin{quote}
Before the new one: say the Sanskrit for a necklace, then the Sanskrit for a crown.
\end{quote}

\subsection*{You'll want to know: \sk{अधोवस्त्रम्}}
\textbf{\sk{अधोवस्त्रम्}} — \emph{adhovastram} — \textquotedblleft{}a lower garment\textquotedblright{}.

\begin{grammarlens}[title={What you've built}]
One more word for this chapter.
\end{grammarlens}

\subsection*{Your turn}
\begin{itemize}
\raggedright
  \item \emph{Say it:} \emph{adhovastram}
  \item \emph{Say it:} \emph{adhovastram}, once more
  \item \emph{Say it:} say \emph{adhovastram}
  \item \emph{Recall:} say the Sanskrit for a necklace, then the Sanskrit for a crown, then say \emph{adhovastram} again
\end{itemize}

\begin{tcolorbox}[breakable,colback=teal!4,colframe=teal!35!black,title={Before you move on}]
What does \sk{अधोवस्त्रम्} mean? (\textquotedblleft{}A lower garment\textquotedblright{}.) Say it once more.
\end{tcolorbox}
\section[\texorpdfstring{\sk{परिधानम्} (paridhānam)}{परिधानम् (paridhānam)}]{\sk{परिधानम्} (paridhānam) — clothing}

\label{lesson:SA-C192-paridhanam}



\begin{quote}
Before the new one: say the Sanskrit for a crown, then the Sanskrit for a lower garment.
\end{quote}

\subsection*{You'll want to know: \sk{परिधानम्}}
\textbf{\sk{परिधानम्}} — \emph{paridhānam} — \textquotedblleft{}clothing\textquotedblright{}.

\begin{grammarlens}[title={What you've built}]
One more word for this chapter.
\end{grammarlens}

\subsection*{Your turn}
\begin{itemize}
\raggedright
  \item \emph{Say it:} \emph{paridhānam}
  \item \emph{Say it:} \emph{paridhānam}, once more
  \item \emph{Say it:} say \emph{paridhānam}
  \item \emph{Recall:} say the Sanskrit for a crown, then the Sanskrit for a lower garment, then say \emph{paridhānam} again
\end{itemize}

\begin{tcolorbox}[breakable,colback=teal!4,colframe=teal!35!black,title={Before you move on}]
What does \sk{परिधानम्} mean? (\textquotedblleft{}Clothing\textquotedblright{}.) Say it once more.
\end{tcolorbox}
\section[\texorpdfstring{\sk{वसनम्} (vasanam)}{वसनम् (vasanam)}]{\sk{वसनम्} (vasanam) — a garment}

\label{lesson:SA-C192-vasanam}



\begin{quote}
Before the new one: say the Sanskrit for a lower garment, then the Sanskrit for clothing.
\end{quote}

\subsection*{You'll want to know: \sk{वसनम्}}
\textbf{\sk{वसनम्}} — \emph{vasanam} — \textquotedblleft{}a garment\textquotedblright{}.

\begin{grammarlens}[title={What you've built}]
One more word for this chapter.
\end{grammarlens}

\subsection*{Your turn}
\begin{itemize}
\raggedright
  \item \emph{Say it:} \emph{vasanam}
  \item \emph{Say it:} \emph{vasanam}, once more
  \item \emph{Say it:} say \emph{vasanam}
  \item \emph{Recall:} say the Sanskrit for a lower garment, then the Sanskrit for clothing, then say \emph{vasanam} again
\end{itemize}

\begin{tcolorbox}[breakable,colback=teal!4,colframe=teal!35!black,title={Before you move on}]
What does \sk{वसनम्} mean? (\textquotedblleft{}A garment\textquotedblright{}.) Say it once more.
\end{tcolorbox}
